\documentclass[12pt,a4paper]{article}
\usepackage{cmap} % Makes the PDF copiable. See http://tex.stackexchange.com/a/64198/25761
\usepackage[T1]{fontenc}
\usepackage[brazil]{babel}
\usepackage[utf8]{inputenc}
\usepackage{amsmath}
\usepackage{amsfonts}
\usepackage{amssymb}
\usepackage{amsthm}
\usepackage{textcomp} % \degree
\usepackage{gensymb} % \degree
\usepackage[usenames,svgnames,dvipsnames]{xcolor}
\usepackage{hyperref}
\usepackage{graphicx}
\usepackage[margin=2cm]{geometry}

\hypersetup{
    colorlinks = true,
    allcolors = {blue}
}

\newcommand*\diff{\mathop{}\!\mathrm{d}}
\newcommand*\sen{\operatorname{sen}}
\newcommand*\tg{\operatorname{tg}}
\newcommand*\abs[1]{\left|#1\right|}

\newcommand*\R{\mathbb{R}}

\newcommand*\tipo{PROVA IV}
\newcommand*\turma{ELE141-01U}
\newcommand*\disciplina{CDI0001}
\newcommand*\eu{Helder G. G. de Lima}
\newcommand*\data{07/12/2015}

\author{\eu}
\title{\tipo - \disciplina}
\date{\data}

\begin{document}
\thispagestyle{empty}
\newgeometry{margin=2cm,bottom=0.5cm}
\begin{center}
\includegraphics[width=9.0cm]{marca}
\noindent\begin{tabular}{l c c r}
  \textbf{\disciplina (\turma)}
& \textbf{\tipo}
& \textbf{\data}
\end{tabular}
\\ Prof. \eu\footnote{
Este é um material de acesso livre distribuído sob os termos da licença \href{https://creativecommons.org/licenses/by-sa/4.0/deed.pt_BR}{Creative Commons BY-SA 4.0}.}
\end{center}

\noindent Nome do(a) aluno(a): \rule{13cm}{0.01cm}

%\section*{Instruções}
\begin{center}\fbox{
\begin{minipage}{14cm}

{\footnotesize
\begin{itemize}
\renewcommand{\theenumi}{\Roman{enumi}}
\item Identifique-se em todas as folhas.
\item Mantenha o celular e os demais equipamentos eletrônicos desligados durante a prova.
\item Justifique cada resposta com cálculos ou argumentos baseados na teoria estudada.
\end{itemize}
}

\end{minipage}
}
\end{center}


%\section*{Questões}

\begin{enumerate}

\item Calcule apenas 5 das 6 integrais a seguir:

\begin{enumerate}
\item (2,0) $\displaystyle \int x^3 \cos(x) \diff x$
\item (2,0) $\displaystyle \int \frac{10 x+14}{x^2+x-20} \diff x$
\item (2,0) $\displaystyle \int \cos(3x)\left[ \cos(2x)+2\sen(x)) \right] \diff x$
\item (2,0) $\displaystyle \int \frac{ \ln(\ln(\ln(x))) }{ x \ln(x) } \diff x$
\item (2,0) $\displaystyle \int \frac{5x-1}{x^2+5} \diff x$
\item (2,0) $\displaystyle \int \tg^6(x) \sec^4(x)\diff x$
\end{enumerate}
\end{enumerate}

\begin{center}
BOA PROVA!!!
\end{center}
\begin{center}
(E BOAS FÉRIAS ;-)
\end{center}

\newpage
\restoregeometry
\section*{Respostas e observações}

\begin{enumerate}

\item \textit{Calcule apenas 5 das 6 integrais a seguir:}

\begin{enumerate}
\item $\displaystyle \int x^3 \cos(x) \diff x$

Considerando $u = x^3$ e $\diff v = \cos(x) \diff x$, tem-se $\diff u=3x^2 \diff x$ e $v = \sen(x)$ e
\begin{align*}
\int x^3 \cos(x) \diff x
= x^3 \sen(x) - \int 3x^2 \sen(x) \diff x
\end{align*}

Tomando $a = 3x^2$ e $\diff b = \sen(x) \diff x$, tem-se $\diff a=6x \diff x$ e $b = -\cos(x)$, de modo que
\begin{align*}
\int x^3 \cos(x) \diff x
& = x^3 \sen(x) - \left( - 3x^2 \cos(x) - \int -6x \cos(x) \diff x \right) \\
& = x^3 \sen(x) + 3x^2 \cos(x) - \int 6x \cos(x) \diff x
\end{align*}

Agora, sendo $p = 6x$ e $\diff q = \cos(x) \diff x$, tem-se $\diff p=6 \diff x$ e $q = \sen(x)$, e então
\begin{align*}
\int x^3 \cos(x) \diff x
& = x^3 \sen(x) + 3x^2 \cos(x) - \left( 6x \sen(x) - \int 6\sen(x) \diff x \right) \\
& = x^3 \sen(x) + 3x^2 \cos(x) - 6x \sen(x) - 6 \cos(x) + c \\
& = (x^3 - 6x) \sen(x) + (3x^2 - 6) \cos(x) + c.
\end{align*}

\item $\displaystyle \int \frac{10 x+14}{x^2+x-20} \diff x$

Observando que $x^2 + x - 20 = (x-4)(x+5)$, a decomposição em frações parciais da função a ser integrada é obtida por
\[
\frac{10 x + 14}{(x-4)(x+5)}
= \frac{A}{x-4} + \frac{B}{x+5}
= \frac{A(x+5) + B(x-4)}{(x-4)(x+5)}
= \frac{(A+B)x +(5A-4B)}{(x-4)(x+5)}.
\]
Comparando os numeradores, resulta que $A$ e $B$ devem ser soluções do sistema
\[
\begin{cases}
A + B = 10 \\
5A - 4B = 14
\end{cases}
\]
Multiplicando a primeira equação por $4$ e somando com a segunda, resulta que $9A = 54$, ou seja, $A=6$. Consequentemente, $B = 10-6 = 4$. Então
\begin{align*}
\int \frac{10 x+14}{x^2+x-20} \diff x
& = \int \frac{6}{x-4} + \frac{4}{x+5} \diff x
  = 6 \int \frac{1}{x-4} \diff x + 4 \int \frac{1}{x+5} \diff x \\
& = 6 \ln\abs{ x-4 } + 4 \ln\abs{x+5} + c
  = \ln\abs{ (x-4)^6 (x+5)^4} + c.
\end{align*}

\item $\displaystyle \int \cos(3x)\left[ \cos(2x)+2\sen(x)) \right] \diff x$

Levando em conta que
\[
\begin{cases}
\cos(3x)\cos(2x) = \frac{1}{2} \left[ \cos(3x-2x) +\cos(3x+2x) \right]
                 = \frac{1}{2} \left[ \cos(x)+\cos(5x) \right] \\
\cos(3x)\sen(x) = \frac{1}{2} \left[ \sen(x-3x) +\sen(x+3x) \right]
                = \frac{1}{2} \left[ \sen(-2x)+\sen(4x) \right]
\end{cases}
\]
tem-se
\begin{align*}
&\int \cos(3x)\left[ \cos(2x)+2\sen(x)) \right] \diff x \\
& = \int \cos(3x)\cos(2x) \diff x + 2 \int \cos(3x)\sen(x)) \diff x \\
& = \int \frac{1}{2} \left[ \cos(x)+\cos(5x) \right] \diff x
+ 2 \int \frac{1}{2} \left[ \sen(-2x)+\sen(4x) \right] \diff x \\
& = \frac{1}{2}\left[ \sen(x) + \frac{\cos(5x)}{5} \right] +
    \left[ \frac{ -\cos(-2x)}{-2} - \frac{\cos(4x)}{4} \right] + c \\
& = \frac{\sen(x)}{2} + \frac{\cos(5x)}{10} +
 \frac{ \cos(2x)}{2} - \frac{\cos(4x)}{4} + c.
\end{align*}

\item $\displaystyle \int \frac{ \ln(\ln(\ln(x))) }{ x \ln(x) } \diff x$

Definindo $w = \ln(\ln(x))$, tem-se $\diff w = \frac{1}{x \ln(x)} \diff x$ e portanto
\[
  \int \frac{ \ln(\ln(\ln(x))) }{ x \ln(x) } \diff x
= \int \ln(\ln(\ln(x))) \frac{1}{ x \ln(x) } \diff x
= \int \ln( w ) \diff w
= \int \ln( w ) \cdot 1 \diff w.
\]
Integrando por partes, com $u = \ln(w)$ e $dv = 1 \diff w$, tem-se $\diff u = \frac{1}{w} \diff w$, $v = w$ e
\[
  \int \ln( w ) \cdot 1 \diff w
= w \ln(w) - \int w \frac{1}{w} \diff w
= w \ln(w) - w + c
= w \cdot [ \ln(w) - 1 ] + c
\]
Assim,
\[
  \int \frac{ \ln(\ln(\ln(x))) }{ x \ln(x) } \diff x
= w ( \ln(w) - 1) + c
= \ln(\ln(x)) \cdot [ \ln( \ln(\ln(x)) ) - 1 ] + c
\]

\item $\displaystyle \int \frac{5x-1}{x^2+5} \diff x$

\begin{align*}
  \int \frac{5x-1}{x^2+5} \diff x
& = 5\int \frac{x}{x^2+5} \diff x
   - \int \frac{1}{x^2+5} \diff x \\
& = \frac{5}{2} \int \frac{2x}{x^2+5} \diff x
   -\frac{1}{5} \int \frac{1}{\left( \frac{x}{\sqrt{5}}\right)^2 + 1} \diff x
\end{align*}

Fazendo $u = x^2+5$ e $v = \frac{x}{\sqrt{5}}$, resulta que $\diff u = 2x \diff x$, $\diff v = \frac{1}{\sqrt{5}} \diff x$ e então
\begin{align*}
  \int \frac{5x-1}{x^2+5} \diff x
= \frac{5}{2} \int \frac{1}{u} \diff u
 -\frac{\sqrt{5}}{5} \int \frac{1}{v^2 + 1} \diff v \\
= \frac{5}{2} \ln \abs{u}
 -\frac{\sqrt{5}}{5} \arctan(v) + c \\
= \frac{5}{2} \ln \abs{x^2+5}
 -\frac{\sqrt{5}}{5} \arctan\left(\frac{x}{\sqrt{5}}\right) + c.
\end{align*}

\item $\displaystyle \int \tg^6(x) \sec^4(x)\diff x$
Usando o fato de que $\sec^2(x) = \tg^2(x) + 1$, tem-se
\[
  \int \tg^6(x) \sec^4(x)\diff x
= \int \tg^6(x) (\tg^2(x) + 1) \sec^2(x)\diff x
= \int (\tg^8(x) + \tg^6(x) ) \sec^2(x)\diff x
\]
Tomando $u = tg(x)$, tem-se $\diff u = \sec^2(x) \diff x$ e
\[
  \int \tg^6(x) \sec^4(x)\diff x
= \int (u^8 + u^6 ) \diff u
= \frac{u^9}{9} + \frac{u^7}{7} + c
= \frac{\tg^9(x)}{9} + \frac{\tg^7(x)}{7} + c.
\]


\end{enumerate}
\end{enumerate}
\end{document}
